% !PRINTED 414
\weberpage{414}
%<<CONT
\gesperrt{IV. If $k$ is relatively prime to $n$ and $r$ runs through the sequence of the primitive $n$th roots of unity, then $r^k$ runs through the same sequence of numbers, though in a different order.}

Finally we cite the further theorem: --

\gesperrt{V. If $n$ is a prime, then every $n$th root of unity, with the exception of 1, is a primitive $n$th root of unity.}

\begin{center}
\S.\,133.\\[0.3em]
\gesperrt{Equations for the primitive roots of unity of degree $n$.}
\end{center}

All $n$th roots of unity are, as we have seen, roots of an equation $f_n(x) = 0$, where
\weq{1}{f_n(x) = x^n - 1}
is understood. If now we free the function $f_n(x)$ from all factors which it has in common with other functions of the same form $f_{n_1}(x)$, which is done by rational operations (\S.\,6), we obtain an equation $X_n = 0$, which is satisfied by the primitive $n$th roots of unity and by these only, and $X_n$ has the form
% [mirrored 2026-08-07] canon ruled the v/nu question for this page: Greek \nu
\weq{2}{X_n = x^\nu + a_1 x^{\nu-1} + \cdots + a_\nu.}
The degree $\nu$ is equal to $\varphi(n)$, and the coefficients $a_1, a_2 \ldots a_\nu$ are \gesperrt{rational numbers}.

In seeking the common factors of $f_n$ and $f_{n_1}$ we may, for $n_1$, restrict ourselves to the divisors of $n$. How the function $X_n$ may be formed in a simple way we shall presently see more closely. But first we prove a general theorem concerning these functions.

The $n$th roots of unity comprise all primitive $\mu$th roots of unity, where $\mu$ is any divisor of $n$, $n$ itself and 1 included, since it is precisely the unity 1 itself that is to be regarded as the primitive first root of unity. If, then, we let $\mu$ run through all divisors of $n$, and observe that two different $X_\mu$ never have a common divisor, and that neither $f_n(x)$ nor $X_\mu$ contains multiple factors, it follows that
\weq{3}{f_n(x) = \Pi X_\mu,}
%CONT>>
